\section{Methods}\label{sec:methods}

\subsection{Prospective Multi-Agent Validation Procedure}
\label{sec:methods_metagent}

No multi-agent simulation result is reported or claimed in this manuscript. The prospective procedure for testing the multi-scale claims, with its three arms (among them the frozen-attention-prior control), predeclared formation event, persistence window, primary statistic, seed count, and pass rules, is stated in Section~\ref{sec:participatory}; it has not been executed.

\subsection{Computational Details}

No numerical result of this manuscript depends on simulation. The analytic gradient formulas of Appendix~\ref{app:gradients} were checked against central finite differences in double precision, with the tolerances stated there. Several appendices (Appendices~\ref{app:gradients}, \ref{app:covariance_dynamics}, \ref{app:conditional_representation}, \ref{app:augmented_joint}, and~\ref{sec:meta_agent_rg}) and the binary-registration example of Section~\ref{sec:binary_registration} report small deterministic numerical witnesses. Each is a direct double-precision evaluation of the stated closed forms or iterations on the stated instance, and none is a simulation of the multi-agent dynamics.

\subsection{Large Language Model Assistance}

Claude Sonnet 4.5 was used as a programming and typesetting assistant. It generated and refactored Python implementations under author direction, drafted and edited prose, and assisted with bibliographic lookup. Claude Opus 5.5 assisted with the revision of this version of the manuscript under author direction, including numerical checks of stated identities. Neither was used to generate experimental results or accepted as an authority for mathematical correctness. Mathematical derivations and citations remain the author's responsibility.

\section{Statements and Declarations}\label{sec:statements}

\subsection*{Competing Interests}
The author declares no competing interests, either financial or non-financial, related to the work submitted for publication.

\subsection*{Author Contributions}
R.C.D. conceived the theoretical framework, performed the mathematical derivations, designed and wrote the multi-agent reference implementation, and wrote the manuscript. Large language model assistance was limited to programming support, fact checking, numerical checking, and drafting and editing prose under author direction, as detailed in the Methods section.

\subsection*{Data Availability}
This manuscript reports no empirical dataset and no simulation output; the validation procedure of Section~\ref{sec:methods_metagent} has not been executed, so no simulation data exist for it.

\subsection*{Code Availability}
A reference implementation of the multi-agent dynamics is publicly available at \url{https://github.com/cdenn016/MAgent_Model}; no argument in this manuscript depends on it.

\subsection*{Funding}
This research received no specific grant from any funding agency in the public, commercial, or not-for-profit sectors.

\begin{acknowledgments}
The author would like to thank Christine King Carter for her generous support, love, and understanding which allowed this work to be pursued.
\end{acknowledgments}

\appendix

\section{\texorpdfstring{Explicit Gradient Expressions for $\mathrm{GL}(K_q)$ Gaussian Agents}{Explicit Gradient Expressions for GL(Kq) Gaussian Agents}}\label{app:gradients}

This appendix collects the gradients and natural gradients of the consensus free-energy functional $\mathcal{F}$ of Section~\ref{sec:free_energy} for Gaussian beliefs with $\mathrm{GL}^+(K_q)$ transport, the connected component containing the transports $\Omega_{ij} = \exp(\phi_i)\exp(-\phi_j)$ of Eq.~\eqref{eq:transport_def}. The compact restriction $\mathrm{SO}(K_q)$ is the special case in which the frame coefficients $\phi_i$ are skew-symmetric, so that each factor $\exp(\phi_i)$ and hence every transport $\Omega_{ij}$ is a rotation. The Gaussian-sector operations are covariant under the congruence action of $\mathrm{GL}(K_q)$ on the covariance fiber; the frame sector is discussed with the update equations.

\subsection*{Setup and Notation}

Agent $i$ maintains primitive belief $q_i = \mathcal{N}(\mu_i, \Sigma_i)$ and primitive generative model $s_i = \mathcal{N}(\mu_{s,i}, \Sigma_{s,i})$, with derived prior $p_i = \mathcal{N}(\mu_{p,i}, \Sigma_{p,i})$ on the state fiber and derived hyper-prior $r_i = \mathcal{N}(\mu_{r,i}, \Sigma_{r,i})$ on the model fiber, supplied by the enveloping meta-agent through the lagged predictive prior of Eq.~\eqref{eq:cross_scale_shadow}. The frame of agent $i$ is the group element $U_i = \exp(\phi_i)$. The state-fiber transport $\Omega_{ij} = U_iU_j^{-1} \in G$ acts via the representation $\rho_{\mathrm{state}}: G \to \mathrm{GL}(K_q)$, and the model-fiber transport $\tilde\Omega_{ij}$ acts via $\rho_{\mathrm{model}}$. This appendix computes the gradients on the state-fiber sector $(q_i, p_i, U_i)$; the model-fiber sector $(s_i, r_i)$ has parallel structure under the substitution $(\beta, \Omega, q, p, \tau) \to (\gamma, \tilde\Omega, s, r, \tau_\gamma)$.

The attention temperatures are $\tau$ on the belief fiber and $\tau_\gamma$ on the model fiber, as in Eq.~\eqref{eq:pointwise_free_energy}; with the parameterization $\tau = \kappa\sqrt{K_q}$ and $\tau_\gamma = \kappa\sqrt{K_m}$, the dimensionless $\kappa > 0$ is a declared constant, not a variational parameter of $\mathcal{F}$. The computation is made at the base point $c_*$ with every agent present, so the receiver-presence weight of each row equals one. Condition numbers are written $\operatorname{cond}(\cdot)$ to avoid a collision with $\kappa$. The pair energies and the attention rows are
\begin{equation}
E_{ij} := D_{\mathrm{KL}}(q_i \| \Omega_{ij}[q_j]),
\qquad
\beta_{ij} = \frac{\pi_{ij}\exp\left(-E_{ij}/\tau\right)}{\sum_m \pi_{im}\exp\left(-E_{im}/\tau\right)},
\end{equation}
so that $\beta_{ij} = \beta_{ij}(\{q\},\{U\})$ depends on the beliefs and frames of the receiver $i$ and of every source $m$ in its row. The scalar $\alpha_i$ is the per-agent prior precision of Eq.~\eqref{eq:free_energy_adaptive}, with optimal value $\alpha_i^* = c_0/(b_0 + D_{\mathrm{KL}}(q_i \| p_i))$ from Eq.~\eqref{eq:state_dependent_alpha}; the constant-precision limit $\alpha_i = 1$ recovers the unit-weight self term of Eq.~\eqref{eq:free_energy_functional_final}. On the model fiber a fixed hyper-prior weight $\lambda_h$ plays the analogous role. A diagonal extension would replace the scalar $\alpha_i$ coordinate-wise by Eq.~\eqref{eq:state_dependent_alpha_per_coord}; it is not developed here.

\subsection*{The Canonical Envelope Gradient}

The dynamics of the main text descend the reduced functional $\mathcal{F}_{\mathrm{red}}$, in which every attention row has been replaced by its Gibbs value $-\tau\log Z_i$ (Section~\ref{sec:envelope_theorem}). At the Gibbs row the derivative of the row block with respect to its weights vanishes, so the gradient of the row value is the gradient taken at frozen weights; this is the attention-reduced potential proposition of~\cite[Ch.~6]{DennisGaugeVFE2026}, whose potential $-\log Z_i$ differs from the row value here by the factor $\tau$. The reduced slow descent is then the natural-gradient flow of the reduced functional by the envelope identity of~\cite[Ch.~8]{DennisGaugeVFE2026}. The adaptive-precision block has the same cancellation (see the discussion following Eq.~\eqref{eq:alpha_product_rule_itfb}), so $\alpha_i$ enters with its optimal value $\alpha_i^*$. A state variable $x_i$ of agent $i$ appears in the row of agent $i$ and, as the sender, in the row of every agent $k$ that attends to $i$. For $x_i \in \{\mu_i, \Sigma_i, U_i\}$ the complete gradient is therefore
\begin{equation}
\nabla_{x_i}\mathcal{F}_{\mathrm{red}} = \alpha_i^*\nabla_{x_i}D_{\mathrm{KL}}(q_i \| p_i) + \sum_j \beta_{ij}\nabla_{x_i}E_{ij} + \sum_k \beta_{ki}\nabla_{x_i}E_{ki} + \nabla_{x_i}\mathcal{F}_{\mathrm{obs},i},
\label{eq:canonical_gradient_general}
\end{equation}
with every weight at its Gibbs value, $\mathcal{F}_{\mathrm{obs},i}$ the observation term of agent $i$, and the model channel contributing terms of the same form, with $(\gamma_{ij}, \tilde\Omega_{ij}, s_i, r_i, \lambda_h)$ in place of $(\beta_{ij}, \Omega_{ij}, q_i, p_i, \alpha_i^*)$, to the gradients with respect to the model variables and, when the model-fiber transport is generated by the same frames, with respect to $U_i$. The observation and self terms carry no frame dependence.

\paragraph{Belief mean and covariance.}
Write $A_{ij} := \Omega_{ij}\Sigma_j\Omega_{ij}^\top$ for the transported covariance of the source. Differentiating the Gaussian relative entropy of Appendix~\ref{app:gaussian_kl}, the self and receiver pieces are
\begin{equation}
\begin{aligned}
\frac{\partial D_{\mathrm{KL}}(q_i \| p_i)}{\partial \mu_i} &= \Sigma_{p,i}^{-1}(\mu_i - \mu_{p,i}),
&\qquad
\frac{\partial D_{\mathrm{KL}}(q_i \| p_i)}{\partial \Sigma_i} &= \frac{1}{2}\left(\Sigma_{p,i}^{-1} - \Sigma_i^{-1}\right), \\
\frac{\partial E_{ij}}{\partial \mu_i} &= A_{ij}^{-1}(\mu_i - \Omega_{ij}\mu_j),
&\qquad
\frac{\partial E_{ij}}{\partial \Sigma_i} &= \frac{1}{2}\left(A_{ij}^{-1} - \Sigma_i^{-1}\right).
\end{aligned}
\end{equation}
Differentiating $E_{ki} = D_{\mathrm{KL}}(q_k \| \Omega_{ki}[q_i])$ with respect to the sender's mean and covariance, using the precision transport of Eq.~\eqref{eq:precision_transport}, gives
\begin{align}
\frac{\partial E_{ki}}{\partial \mu_i} &= \Lambda_{q_i}\Omega_{ki}^{-1}\left(\Omega_{ki}\mu_i - \mu_k\right), \\
\frac{\partial E_{ki}}{\partial \Sigma_i} &= \frac{1}{2}\Omega_{ki}^{\top}\left[\tilde{\Lambda}_{q_i}^{(k)} - \tilde{\Lambda}_{q_i}^{(k)}\left(\Sigma_k + d_{ki}d_{ki}^{\top}\right)\tilde{\Lambda}_{q_i}^{(k)}\right]\Omega_{ki},
\end{align}
where $\Lambda_{q_i} = \Sigma_i^{-1}$, $\tilde{\Lambda}_{q_i}^{(k)} := (\Omega_{ki}\Sigma_i\Omega_{ki}^{\top})^{-1} = \Omega_{ki}^{-\top}\Lambda_{q_i}\Omega_{ki}^{-1}$ is agent $i$'s precision transported into agent $k$'s frame, and $d_{ki} := \Omega_{ki}\mu_i - \mu_k$ is the mean residual in agent $k$'s frame. The rank-one piece $-\tilde{\Lambda}_{q_i}^{(k)}d_{ki}d_{ki}^{\top}\tilde{\Lambda}_{q_i}^{(k)}$ vanishes at consensus and contributes a stiffness aligned with $d_{ki}$ away from it. The inverse-transpose form of Eq.~\eqref{eq:precision_transport} is what makes these expressions valid for general $\Omega_{ki} \in \mathrm{GL}^{+}(K_q)$, where $\Omega_{ki}^{-1}$ does not reduce to $\Omega_{ki}^{\top}$. Substituting into Eq.~\eqref{eq:canonical_gradient_general} gives the complete mean and covariance gradients,
\begin{equation}
\begin{aligned}
\nabla_{\mu_i}\mathcal{F}_{\mathrm{red}} &= \alpha_i^*\Sigma_{p,i}^{-1}(\mu_i - \mu_{p,i}) + \sum_j \beta_{ij}A_{ij}^{-1}(\mu_i - \Omega_{ij}\mu_j) \\
&\quad + \sum_k \beta_{ki}\Lambda_{q_i}\Omega_{ki}^{-1}\left(\Omega_{ki}\mu_i - \mu_k\right) + \nabla_{\mu_i}\mathcal{F}_{\mathrm{obs},i},
\end{aligned}
\label{eq:canonical_gradient_mu}
\end{equation}
\begin{equation}
\begin{aligned}
\nabla_{\Sigma_i}\mathcal{F}_{\mathrm{red}} &= \frac{\alpha_i^*}{2}\left(\Sigma_{p,i}^{-1} - \Sigma_i^{-1}\right) + \frac{1}{2}\sum_j \beta_{ij}\left(A_{ij}^{-1} - \Sigma_i^{-1}\right) \\
&\quad + \frac{1}{2}\sum_k \beta_{ki}\Omega_{ki}^{\top}\left[\tilde{\Lambda}_{q_i}^{(k)} - \tilde{\Lambda}_{q_i}^{(k)}\left(\Sigma_k + d_{ki}d_{ki}^{\top}\right)\tilde{\Lambda}_{q_i}^{(k)}\right]\Omega_{ki} + \nabla_{\Sigma_i}\mathcal{F}_{\mathrm{obs},i},
\end{aligned}
\label{eq:canonical_gradient_sigma}
\end{equation}
plus the model-channel terms. Eq.~\eqref{eq:canonical_gradient_mu} is Eq.~\eqref{eq:total_gradient_mu} with the sender index relabeled $k \to j$, and Eq.~\eqref{eq:total_gradient_mu} gives the observation terms for a Gaussian identity readout.

\paragraph{Frame sector: the transported relative entropy.}
The frames enter only through the links. For Gaussians write $A := \Omega_{ij}\Sigma_j\Omega_{ij}^\top$, $r := \mu_i - \Omega_{ij}\mu_j$, and $\Lambda := A^{-1}$. The Gaussian relative entropy has three $\Omega_{ij}$-dependent pieces ($\log|A|$, $\operatorname{tr}(\Lambda\Sigma_i)$, $r^\top\Lambda r$), and the group is $\mathrm{GL}^+(K_q)$, so the directional derivative lives in the full Lie algebra $\mathfrak{gl}(K_q)$ with no skew-symmetry constraint. Along $X \in \mathfrak{gl}(K_q)$ acting by $\Omega_{ij} \to e^{tX}\Omega_{ij}$, with $\partial_X A = XA + AX^\top$ and $\partial_X r = -X\Omega_{ij}\mu_j$, the derivative is
\begin{equation}
\partial_X E_{ij} = \operatorname{tr}\Big[X^\top\Big(I - \Lambda\Sigma_i - \Lambda r(\Omega_{ij}\mu_j)^\top - \Lambda r r^\top\Big)\Big],
\label{eq:kl_omega_directional}
\end{equation}
where the leading identity is the $\log\det$ contribution $\partial_X\log|A| = 2\operatorname{tr}(X)$ carried at the half weight of the Gaussian relative entropy, and the quadratic term uses $\Lambda A = I$ to reduce $\Lambda r r^\top\Lambda A$ to $\Lambda r r^\top$. Writing $N_{ij} := I - \Lambda\Sigma_i - \Lambda r(\Omega_{ij}\mu_j)^\top - \Lambda r r^\top$ for the bracket of Eq.~\eqref{eq:kl_omega_directional}, the unconstrained matrix derivative is
\begin{equation}
\frac{\partial E_{ij}}{\partial \Omega_{ij}} = N_{ij}\Omega_{ij}^{-\top} = \Omega_{ij}^{-\top} - \Lambda \Sigma_i \Omega_{ij}^{-\top} - \Lambda(\mu_i - \Omega_{ij}\mu_j)\mu_j^\top - \Lambda r r^\top \Lambda \Omega_{ij}\Sigma_j,
\label{eq:kl_omega_unconstrained}
\end{equation}
where the mean term carries a negative sign because $r$ depends on $\Omega_{ij}$ through $-\Omega_{ij}\mu_j$, and the simplifications $(\Omega_{ij}\mu_j)^\top\Omega_{ij}^{-\top} = \mu_j^\top$ and $\Omega_{ij}^{-\top} = \Lambda\Omega_{ij}\Sigma_j$ reduce the remaining factors. At consensus, $q_i = \Omega_{ij}[q_j]$, the bracket $N_{ij}$ vanishes. The compact subgroup $\mathrm{SO}(K_q)$ is the special case $X \in \mathfrak{so}(K_q) \subset \mathfrak{gl}(K_q)$ ($X^\top = -X$), in which the $\log\det$ piece vanishes because $\operatorname{tr}(X) = 0$ and the gradient is obtained by projecting onto $\mathfrak{so}(K_q)$ through antisymmetrization.

\paragraph{Frame sector: gradient in the group variable.}
Since $\Omega_{ij} = U_iU_j^{-1}$, a variation of $U_i$ moves the receiver link by $\delta\Omega_{ij} = (\delta U_i)U_j^{-1}$ and the sender link by $\delta\Omega_{ki} = -\Omega_{ki}(\delta U_i)U_i^{-1}$. Combining with Eq.~\eqref{eq:kl_omega_unconstrained} and the identities $\Omega_{ij}^{-\top}U_j^{-\top} = U_i^{-\top}$ and $\Omega_{ki}^{-\top}U_i^{-\top} = U_k^{-\top}$ gives
\begin{equation}
\nabla_{U_i}E_{ij} = N_{ij}U_i^{-\top},
\qquad
\nabla_{U_i}E_{ki} = -\Omega_{ki}^{\top}N_{ki}U_k^{-\top},
\end{equation}
with $N_{ki}$ the bracket of Eq.~\eqref{eq:kl_omega_directional} for the pair $(k, i)$. The complete frame gradient is therefore
\begin{equation}
\nabla_{U_i}\mathcal{F}_{\mathrm{red}} = \sum_j \beta_{ij}N_{ij}U_i^{-\top} - \sum_k \beta_{ki}\Omega_{ki}^{\top}N_{ki}U_k^{-\top},
\label{eq:canonical_gradient_frame}
\end{equation}
plus the model-channel terms of the same form when the model-fiber transport $\tilde\Omega_{ij}$ is generated by the same frames through $\rho_{\mathrm{model}}$ (tied frames); with independent model-fiber frames the model-channel gradient is taken with respect to those frames instead. Every term vanishes at consensus. In the frame of the receiving agent of each row, the body covector $U_i^\top\nabla_{U_i}\mathcal{F}_{\mathrm{red}}$ reads $\sum_j\beta_{ij}U_i^\top N_{ij}U_i^{-\top} - \sum_k\beta_{ki}U_k^\top N_{ki}U_k^{-\top}$, so each pair contributes its bracket conjugated into the frame of the agent whose row contains it.

\paragraph{Frame sector: chart gradients.}
In the chart $U_i = \exp(\phi_i)$ with generators $G_a$, the chain rule gives $\partial\mathcal{F}_{\mathrm{red}}/\partial\phi_i^a = \operatorname{tr}[(D\exp_{\phi_i}[G_a])^\top\nabla_{U_i}\mathcal{F}_{\mathrm{red}}]$. For the receiver link the exact Fr\'echet derivative of the matrix exponential is
\begin{equation}
\frac{\partial \Omega_{ij}}{\partial \phi_i^a} = \left(\int_0^1 e^{(1-s)\phi_i} G_a e^{s\phi_i} \mathrm{d}s\right) e^{-\phi_j},
\label{eq:frechet_dexp}
\end{equation}
and for the sender link, where $\phi_i$ enters through the right factor $\exp(-\phi_i)$,
\begin{equation}
\frac{\partial \Omega_{ki}}{\partial \phi_i^a} = -\exp(\phi_k)J_a,
\qquad
J_a := \int_0^1 e^{-(1-s)\phi_i} G_a e^{-s\phi_i} \mathrm{d}s,
\label{eq:frechet_dexp_sender}
\end{equation}
so that
\begin{equation}
\frac{\partial E_{ij}}{\partial \phi_i^a} = \operatorname{tr}\left[\left(\frac{\partial \Omega_{ij}}{\partial \phi_i^a}\right)^\top N_{ij}\Omega_{ij}^{-\top}\right],
\qquad
\frac{\partial E_{ki}}{\partial \phi_i^a} = -\operatorname{tr}\left[J_a^\top (e^{\phi_k})^\top N_{ki}\Omega_{ki}^{-\top}\right].
\end{equation}
When $G_a$ commutes with $\phi_i$ the integrals collapse, $\partial \Omega_{ij}/\partial \phi_i^a = G_a \Omega_{ij}$ and $\partial \Omega_{ki}/\partial \phi_i^a = -\Omega_{ki} G_a$, and the sender derivative becomes
\begin{equation}
\frac{\partial E_{ki}}{\partial \phi_i^a} = -\operatorname{tr}\left[G_a^\top\left(\Omega_{ki}^\top N_{ki} \Omega_{ki}^{-\top}\right)\right] \qquad \text{(abelian approximation)},
\label{eq:sender_phi_gradient}
\end{equation}
while the receiver derivative becomes the projection of Eq.~\eqref{eq:kl_omega_directional} along $X = G_a$. Used as an approximation for noncommuting frames, the abelian form is local and does not rely on compactness of $G$ (the group $\mathrm{GL}^+(K_q)$ is noncompact); its relative error in $\partial\Omega/\partial\phi$ is about $0.37\|\phi_i\|_F$ for small frames, reaching 1 percent at $\|\phi_i\|_F \approx 0.03$, 10 percent at $\|\phi_i\|_F \approx 0.3$, and about 40 percent at $\|\phi_i\|_F \approx 1.4$ (median over generic frames in $\mathfrak{gl}(3)$). The group-variable form Eq.~\eqref{eq:canonical_gradient_frame} requires no Fr\'echet integral and is exact at every frame.

\subsection*{Natural Gradients and Frame Covariance}

The natural gradient raises the index of the gradient with the Fisher metric of Appendix~\ref{app:fisher_gaussian}: $\tilde{\nabla}_{\mu_i} \mathcal{F} = \Sigma_i \nabla_{\mu_i} \mathcal{F}$ on the mean and $\tilde{\nabla}_{\Sigma_i} \mathcal{F} = 2\Sigma_i (\nabla_{\Sigma_i} \mathcal{F}) \Sigma_i$ on $\mathcal{S}_{++}^{K_q}$.

\textbf{Frame covariance.} Under the congruence action of $g \in \mathrm{GL}(K_q)$, $\Sigma_i \mapsto g\Sigma_i g^\top$, which preserves the cone $\mathcal{S}_{++}^{K_q}$ for any invertible $g$, the Euclidean gradient $\nabla_{\Sigma_i}\mathcal{F}$ is a covector on the symmetric matrices and transforms by the inverse-transpose congruence, whereas the natural gradient transforms like $\Sigma_i$ itself:
\begin{align}
\Sigma_i &\mapsto g\Sigma_i g^\top, \\
\nabla_{\Sigma_i} \mathcal{F} &\mapsto g^{-\top}(\nabla_{\Sigma_i} \mathcal{F})g^{-1}, \\
\tilde{\nabla}_{\Sigma_i} \mathcal{F} &\mapsto g(\tilde{\nabla}_{\Sigma_i} \mathcal{F})g^\top.
\end{align}
The two laws coincide if and only if $g$ is orthogonal, $g^{-\top} = g$, that is $g \in \mathrm{O}(K_q)$, so for general $g$ only the natural gradient is covariant on $\mathcal{S}_{++}^{K_q}$; the analogous statement for a constant-metric Gaussian mean flow, that the Euclidean metric is equivariant only under orthogonal reframings, is proved in~\cite[Ch.~G8]{DennisGaugeVFESuppG2026}. Riemannian metrics on $\mathcal{S}_{++}^{K_q}$ invariant under the full congruence action form the two-parameter family $a\operatorname{tr}(\Sigma^{-1}U\Sigma^{-1}V) + b\operatorname{tr}(\Sigma^{-1}U)\operatorname{tr}(\Sigma^{-1}V)$ with $a > 0$ and $a + K_q b > 0$~\cite{Pennec2009}. The Gaussian Fisher metric on the covariance block is the member $a = \tfrac{1}{2}$, $b = 0$~\cite[Ch.~G5]{DennisGaugeVFESuppG2026}, and $2\Sigma(\cdot)\Sigma$ is its gradient structure; the member $a = 1$, $b = 0$ would give $\Sigma(\cdot)\Sigma$. On the unit-determinant directions $\operatorname{tr}(\Sigma^{-1}U) = 0$, which form the $\mathfrak{sl}(K_q)$ sector after whitening, the second term vanishes and all members coincide up to scale. The gradient structure $2\Sigma(\cdot)\Sigma$ is therefore one invariant choice among many, singled out by the Fisher metric.

\subsection*{Exponential-Map Update on the SPD Manifold}

After computing $\tilde{\nabla}_{\Sigma_i}\mathcal{F}$, the covariance is updated through a retraction, a smooth map from the tangent bundle to the manifold whose differential at the zero vector is the identity~\cite{Absil2008}, so that $\Sigma_i$ remains positive definite:
\begin{equation}
\Sigma_i^{\mathrm{new}} = \mathcal{R}_{\Sigma_i}(-\eta_\Sigma \tilde{\nabla}_{\Sigma_i}\mathcal{F}).
\end{equation}
The retraction used is the affine-invariant exponential map,
\begin{equation}
\mathcal{R}_{\Sigma}(\Delta) = \Sigma^{1/2} \exp\left(\Sigma^{-1/2} \Delta \Sigma^{-1/2}\right) \Sigma^{1/2},
\end{equation}
where $\exp(\cdot)$ is the matrix exponential. It preserves positive-definiteness, since the exponential of a symmetric matrix is SPD and congruence preserves the cone, and it is equivariant under congruence, $\mathcal{R}_{A\Sigma A^\top}(A\Delta A^\top) = A \mathcal{R}_{\Sigma}(\Delta) A^\top$ for any invertible $A$. With the eigendecomposition $\Sigma = V D V^\top$,
\begin{equation}
\mathcal{R}_{\Sigma}(\Delta) = V D^{1/2} \exp\left(D^{-1/2} V^\top \Delta V D^{-1/2}\right) D^{1/2} V^\top.
\end{equation}
This is the Riemannian exponential map of the affine-invariant metric, so the update follows the geodesic from $\Sigma$ in direction $\Delta$ for unit parameter.

\subsection*{The Entropy-Suppressed Surrogate}

Replacing each row block by $S_k := \sum_m \beta_{km}E_{km}$, with the Gibbs weights attached but the attention entropy removed, defines a different scalar $S_{\mathrm{surr}}$. Applying the covariance identity of Section~\ref{sec:complete_free_energy} to every row in which $x_i$ appears gives
\begin{equation}
\nabla_{x_i}S_{\mathrm{surr}} = \nabla_{x_i}\mathcal{F}_{\mathrm{red}} - \frac{1}{\tau}\sum_{k}\operatorname{Cov}_{\beta_k}\left(E_{k\cdot}, \nabla_{x_i}E_{k\cdot}\right),
\label{eq:surrogate_covariance_gap}
\end{equation}
where the sum runs over the receiver row $k = i$ and every sender row containing $i$. In a sender row only $E_{ki}$ depends on $x_i$, and the covariance reduces to $\beta_{ki}(E_{ki} - \langle E\rangle_{\beta_k})\nabla_{x_i}E_{ki}$. In the receiver row the surrogate contribution is $\sum_j \beta_{ij}[1 - (E_{ij} - \langle E\rangle_{\beta_i})/\tau]\nabla_{x_i}E_{ij}$. It reduces to the frozen-weight gradient whenever $\operatorname{Cov}_{\beta_i}(E, \nabla E) = 0$, for example when all $E_{ij}$ are equal; it tends to the $\pi$-weighted mean gradient as $\tau \to \infty$; and it tends to $\nabla E_{ij^*}$ as $\tau \to 0$ when the minimizer $j^*$ is unique, since for an energy gap $\Delta > 0$ the response term is bounded by a multiple of $(\Delta/\tau)e^{-\Delta/\tau}$. The response is largest near ties, where energy differences are of order $\tau$. The split of the same quantity into a direct piece weighted by $(1 - E_{ij}/\tau)$ and a reweighting piece $\tau^{-1}\langle E\rangle_{\beta_i}\sum_m\beta_{im}\nabla E_{im}$ is algebraically equal but not centered: a common shift of all energies moves weight between the two pieces, each grows like $E/\tau$ as $\tau \to 0$ while their sum stays bounded, and evaluating them separately cancels catastrophically at small $\tau$. The centered form above avoids this. On the model fiber the same statements hold with $(\gamma, \tilde\Omega, s, \tau_\gamma)$. For the precision weight, differentiating $\alpha_i^*(\mu_i,\Sigma_i)D_{\mathrm{KL}}(q_i\|p_i)$ without the regularizer of Eq.~\eqref{eq:free_energy_adaptive} gives the prefactor $(\alpha_i^*)^2 b_0/c_0$ in place of $\alpha_i^*$ (Eq.~\eqref{eq:alpha_product_rule_itfb}).

Discarding the sender rows defines a further object, the receiver-only field $V_i^{\mathrm{rec}}(x) := \partial_{x_i}S_{i}(x)$, obtained by dropping $\sum_{k \ne i}\partial_{x_i}S_{k}$. It is not the gradient of $S_{\mathrm{surr}}$ and is generically nonintegrable; the receiver and sender terms together form $\nabla S_{\mathrm{surr}}$, whose Hessian is symmetric. Neither the surrogate nor the receiver-only field is the update of the main text, which descends $\mathcal{F}_{\mathrm{red}}$ with the gradients of Eqs.~\eqref{eq:canonical_gradient_general}--\eqref{eq:canonical_gradient_frame}.

\subsection*{Complete Update Equations}

The update uses a Fisher-Rao natural-gradient step on the Gaussian sector and a left-trivialized group retraction on the frame sector (cf.\ Eqs.~\eqref{eq:gauge_natural_gradient}--\eqref{eq:gauge_group_retraction}). Every derivative below is the complete receiver-plus-sender envelope gradient of $\mathcal{F}_{\mathrm{red}}$:
\begin{align}
\mu_i^{t+1} &= \mu_i^t - \eta_\mu \Sigma_i^t \nabla_{\mu_i} \mathcal{F}_{\mathrm{red}}, \\
\Sigma_i^{t+1} &= \left(\Sigma_i^t\right)^{1/2} \exp\left[\left(\Sigma_i^t\right)^{-1/2}\left(-\eta_\Sigma \tilde{\nabla}_{\Sigma_i} \mathcal{F}_{\mathrm{red}}\right)\left(\Sigma_i^t\right)^{-1/2}\right] \left(\Sigma_i^t\right)^{1/2}, \\
\xi_i^t &= \left(U_i^t\right)^{-1}\nabla_{U_i} \mathcal{F}_{\mathrm{red}}, \\
U_i^{t+1} &= U_i^t \exp\left(-\eta_\phi \xi_i^t\right).
\end{align}
Here $\xi_i$ is the left-trivialized frame velocity of Eq.~\eqref{eq:gauge_natural_gradient_def}: the chart natural-gradient vector $\mathcal{G}_{\mathrm{pull}}^{-1}\nabla_{\phi_i}\mathcal{F}$ pushes forward to $\nabla_{U_i}\mathcal{F}$, so its left-trivialized velocity is exactly $U_i^{-1}\nabla_{U_i}\mathcal{F}$ when the generators span $\mathfrak{gl}(K_q)$, and its tangent projection on a subalgebra (Section~\ref{sec:timescales}), with $\nabla_{U_i}\mathcal{F}_{\mathrm{red}}$ given by Eq.~\eqref{eq:canonical_gradient_frame}. Evaluated in this form the velocity costs one $K_q \times K_q$ solve per agent; forming and inverting the $K_q^2 \times K_q^2$ pullback metric explicitly would cost $O(K_q^6)$ per agent and square its condition number. The alternative left-invariant Frobenius preconditioner replaces $\xi_i$ by the body covector $U_i^\top\nabla_{U_i}\mathcal{F}_{\mathrm{red}}$. The body-covector update is equivariant under per-agent left translations with the beliefs co-transformed, and under right translations only by orthogonal $g$; the ambient-pullback update above is equivariant only under orthogonal translations. Neither descends to the quotient by the global right redundancy described below. A right-trivialized update would descend to that quotient at the cost of left equivariance, and no Riemannian metric on $\mathrm{GL}^+(K_q)$ is invariant under both translations~\cite{Milnor1976}. The exponential map in the last line keeps $U_i \in G$ exactly. The learning rates follow the separation declared in Section~\ref{sec:timescales}, with $\eta_\mu \equiv \eta_q$ the belief-fiber rate written for the separately updated mean block.

\subsection*{Computational Complexity}

Per gradient step with $N$ agents and dimension $K_q$, forming the links costs $O(N K_q^3)$ for the frames and their inverses plus $O(N^2 K_q^3)$ for the pairwise products; the transported relative entropies require factorizations of the $O(N^2)$ transported covariances, $O(N^2 K_q^3)$ in total; and the natural-gradient steps and exponential-map updates, with the frame velocity evaluated as $U_i^{-1}\nabla_{U_i}\mathcal{F}_{\mathrm{red}}$, cost $O(K_q^3)$ per agent, $O(N K_q^3)$ in total. The overall complexity is $O(N^2 K_q^3)$ per step, dominated by the pairwise coupling.

\subsection*{Numerical Stability and Verification}

\textbf{Positive-definiteness in finite precision.} In exact arithmetic the exponential-map update maps into $\mathcal{S}_{++}^{K_q}$ for any symmetric $\Delta$, so $\Sigma_i \succ 0$ is preserved by construction. In single precision, with unit roundoff $u = 2^{-24} \approx 5.96 \times 10^{-8}$, a symmetric matrix whose condition number exceeds about $u^{-1} \approx 1.68 \times 10^{7}$ can already be indefinite as stored, before any eigendecomposition is computed; an eigenvalue floor on the stored matrices, or storage of a factor, is what restores definiteness, and a higher-precision eigendecomposition does not. The same bound applies to the transported covariances $A_{ij} = \Omega_{ij}\Sigma_j\Omega_{ij}^\top$ as well as to $\Sigma_i$, and the bound
\begin{equation}
\operatorname{cond}(A_{ij}) \le \operatorname{cond}(\Omega_{ij})^2\operatorname{cond}(\Sigma_j)
\end{equation}
holds, where $\operatorname{cond}(\Omega_{ij})$ grows exponentially with the frame norms, so a bound on the frame norms is needed alongside the eigenvalue floor. Cholesky factorization of an $n \times n$ matrix is backward stable independently of conditioning: the computed factor satisfies $R^\top R = \Sigma + \delta\Sigma$ with $|\delta\Sigma| \le \gamma_{n+1}|R^\top||R|$ componentwise, $\gamma_{n+1} = (n+1)u/(1 - (n+1)u)$~\cite[Thm.~10.3]{Higham2002}. What depends on conditioning is completion and forward error. The factorization runs to completion if $\lambda_{\min}(H) > n\gamma_{n+1}/(1 - \gamma_{n+1})$, where $H = D^{-1}\Sigma D^{-1}$ with $D = \operatorname{diag}(\Sigma)^{1/2}$ is the unit-diagonal scaling of $\Sigma$~\cite[Thm.~10.7]{Higham2002}; since $\lambda_{\max}(H) \ge 1$, this holds in particular when $\operatorname{cond}_2(H)$ is below about $1/(n(n+1)u)$, which for $n = K_q = 10$ in single precision is about $1.5 \times 10^5$.

\textbf{Right-translation redundancy.} Under the global right translation $U_i \mapsto U_i g$ by a single $g \in G$ shared by all agents (Section~\ref{sec:gauge_covariance_principle}), every link $\Omega_{ij} = U_iU_j^{-1} \mapsto U_i g g^{-1} U_j^{-1} = \Omega_{ij}$ is unchanged. The free energy and the belief updates, which depend on the frames only through the links, are therefore unchanged. The frame update is not: since $\nabla_{U_ig}\mathcal{F} = (\nabla_{U_i}\mathcal{F})g^{-\top}$, the velocity becomes $g^{-1}\xi_i g^{-\top}$, and $U_ig\exp(-\eta_\phi g^{-1}\xi_ig^{-\top})$ agrees with $U_i\exp(-\eta_\phi\xi_i)g$ for generic $\xi_i$ only when $g^{-\top} = g$, that is for orthogonal $g$; the body-covector update behaves in the same way. Initial frames related by a nonorthogonal right translation, which describe the same links, therefore evolve to different links. In the chart, $U_ig$ corresponds to $\log(e^{\phi_i}e^{\zeta}) = \phi_i + \psi(\mathrm{ad}_{\phi_i})(\zeta) + O(\|\zeta\|^2)$ with $\zeta = \log g$ and $\psi(z) = z/(1 - e^{-z}) = 1 + z/2 + z^2/12 + \cdots$. The additive shift $\phi_i \mapsto \phi_i + \zeta$ is therefore correct only to first order in both $\phi_i$ and $\zeta$, or exactly in the abelian sector, and the group-level statement $U_i \mapsto U_i g$ is the one to use for the nonabelian redundancy.

\textbf{Gradient verification.} Every formula of this appendix was checked against central finite differences with step $10^{-6}$ in double precision, for $K_q \in \{2, 3\}$, three agents, and generic $\mathrm{GL}^+(K_q)$ frames with $\|\phi_i\|_F \approx 1.5$. The receiver, sender, and self gradients of the mean and covariance, the matrix derivative Eq.~\eqref{eq:kl_omega_unconstrained}, the Fr\'echet integrals Eqs.~\eqref{eq:frechet_dexp} and~\eqref{eq:frechet_dexp_sender}, the exact chart and group-variable frame gradients, the complete envelope gradients, and the surrogate gradients with both precision conventions agree with the finite differences to relative error at most $10^{-7}$; the directional derivative Eq.~\eqref{eq:kl_omega_directional} agrees to $2.3 \times 10^{-7}$. The abelian approximation carries the error quantified above, not machine-precision error.

\section{Intuitive Examples and Framework Extensions}
\label{app:examples}

The derivations of Section~\ref{sec:rock_derivations} are exact consequences of the Gaussian consensus functional under stated limits (status D); their reading as statements about rocks and observers has status S (Table~\ref{tab:claim_status}). None of the quantities is identified with relativistic proper time, inertial mass, or the quantum-mechanical measurement formalism.

\subsection{Macroscopic Objects as Meta-Agents: The Rock Example}

To illustrate how the framework describes familiar objects, consider a rock. In the formalism a rock is not a primitive entity but a meta-agent with hierarchical structure.

A typical rock comprises approximately $10^{23}$ atomic constituents. Each atom is read as a level-0 agent maintaining beliefs about its local environment, primarily the positions and momenta of neighboring atoms. Atoms in a solid are strongly coupled to their neighbors through electromagnetic interactions, which the framework represents by large attention weights $\beta_{\text{atom,neighbor}} \approx 1$. This strong coupling drives belief alignment: atoms in a crystalline lattice maintain highly coherent beliefs about their relative spatial configuration.

When these atomic agents achieve sufficient coherence (small divergences between transported beliefs, in the sense of the threshold detector of Section~\ref{sec:meta_agent_threshold}, which is a proxy for the licensing criterion of Eq.~\eqref{eq:meta_agent_FE_criterion}), they form a level-1 meta-agent, the rock itself. This meta-agent carries a collective belief $q_{\text{rock}}$ representing the consensus configuration of its constituents, a collective prior $p_{\text{rock}}$ encoding the stable lattice structure of the material, and a collective frame $U_{\text{rock}}$, which by Proposition~\ref{prop:parent_frame_gauge} is a declared convention, not an observable. The rock admits an induced metric $G_{\text{rock}}$ by pullback of its collective belief section (Section~\ref{sec:pullback}); the commitment to universal agency of Section~\ref{sec:cognitive_first} does not ascribe phenomenal properties to this geometry.

The rock's Fisher information is large, and the derivation below shows that the pull of a low-precision neighbor on its mean is suppressed by the rock's own small covariance. Its relaxation rate toward its own prior is $\eta_q$, independent of its precision, because the Fisher metric and the stiffness of the prior term scale together; the insensitivity is a property of the coupling, not slow dynamics. In the reading of Section~\ref{sec:mass}, large precision is large stiffness. Near its stationary state the rock's Fisher clock $\vartheta_{\text{rock}}$ (Section~\ref{sec:fisher_arc_length}) advances little, because its belief changes little.

When a human observes the rock, the coupling $\beta_{\text{human,rock}}$ becomes nonzero. The coupling is not direct but is mediated by photons, electromagnetic quanta that couple strongly to both rock atoms (scattering) and human photoreceptors (absorption). The photons carry information about the rock's configuration to the observer, and the human agent's beliefs $q_{\text{human}}$ update toward the transported rock beliefs $\Omega_{\text{human,rock}}[q_{\text{rock}}]$ through free-energy minimization. Within the framework, ``seeing the rock'' is read as agent-agent information coupling mediated by intermediate channels.

Observation is bidirectional: the rock's state also updates in response to the photons that scattered from it. The derivation below shows that its belief mean barely moves, while coupling to any agent lowers its stationary precision by a factor that does not tend to one in the high-precision limit.

\subsubsection{Derivations: insensitivity to low-precision neighbors, cross-observer consensus, and back-action}
\label{sec:rock_derivations}

The statements of the previous paragraphs follow from the natural-gradient flow on $\mathcal{F}$ together with the Gaussian relative-entropy gradients. Three derivations make this explicit.

\paragraph{Insensitivity to low-precision neighbors.} Consider the single-agent free energy
\begin{equation}
\mathcal{F}_i = D_{\mathrm{KL}}(q_i \| p_i) + \beta_{ij} D_{\mathrm{KL}}(q_i \| \Omega_{ij} q_j),
\end{equation}
with Gaussian $q_i = \mathcal{N}(\mu_i, \Sigma_i)$, $p_i = \mathcal{N}(\bar{\mu}_i, \bar{\Sigma}_i)$, and transported neighbor $\Omega_{ij} q_j = \mathcal{N}(\tilde{\mu}_j, \tilde{\Sigma}_{q_j})$. Differentiating the Gaussian relative entropy of Appendix~\ref{app:gaussian_kl} with respect to $\mu_i$ gives
\begin{equation}
\nabla_{\mu_i}\mathcal{F}_i = \bar{\Lambda}_{p_i}(\mu_i - \bar{\mu}_i) + \beta_{ij}\tilde{\Lambda}_{q_j}(\mu_i - \tilde{\mu}_j),
\label{eq:rock_grad_mu}
\end{equation}
where $\bar{\Lambda}_{p_i} = \bar{\Sigma}_i^{-1}$ and $\tilde{\Lambda}_{q_j} = \tilde{\Sigma}_{q_j}^{-1}$. The natural-gradient flow with Fisher preconditioner $g_{\mu\mu} = \Sigma_i^{-1}$ (Appendix~\ref{app:fisher_gaussian}) gives the mean update
\begin{equation}
\dot{\mu}_i = -\eta_q \Sigma_i \nabla_{\mu_i}\mathcal{F}_i = -\eta_q\Sigma_i\left[\bar{\Lambda}_{p_i}(\mu_i - \bar{\mu}_i) + \beta_{ij}\tilde{\Lambda}_{q_j}(\mu_i - \tilde{\mu}_j)\right].
\label{eq:rock_natural_grad_step}
\end{equation}
For a rock with $\Sigma_i = \bar{\Sigma}_i = \delta I$, where $\delta \ll 1$ is the dimensionless precision scale of Section~\ref{sec:mass}, to which no physical units are assigned ($\bar{\Lambda}_{p_i} = \delta^{-1} I$), and a low-precision observer agent $j$ with $\tilde{\Lambda}_{q_j} = O(1)$, substitution yields
\begin{equation}
\dot{\mu}_i = -\eta_q(\mu_i - \bar{\mu}_i) - \eta_q \delta \beta_{ij} \tilde{\Lambda}_{q_j} (\mu_i - \tilde{\mu}_j) = -\eta_q(\mu_i - \bar{\mu}_i) + O(\delta).
\label{eq:rock_inertia}
\end{equation}
The prior-anchor term is $O(1)$ and survives; the coupling to the low-precision observer is $O(\delta)$ and vanishes as $\delta \to 0$. The rock's mean is pinned to its own sharp prior at leading order, while its relaxation rate toward that prior remains $\eta_q$: the Fisher preconditioner $\Sigma_i$ in Eq.~\eqref{eq:rock_natural_grad_step} converts the large precision of the rock into a small response to coupling terms whose own precision is of order one~\cite{Amari2016}.

\paragraph{Cross-observer consensus via photon-mediated coupling.} Consider three kinds of agent: the rock $R$ with precision $\Lambda_R = \delta^{-1}$, a photon channel $\mathrm{ph}$, and human observers $h_k$. Under the scattering event the channel couples to $R$ with weight $\beta_{\mathrm{ph},R}\tilde{\Lambda}_{q_R} \gg \bar{\Lambda}_{\mathrm{ph}}$; setting $\dot{\mu}_{\mathrm{ph}} = 0$ in Eq.~\eqref{eq:rock_grad_mu} gives
\begin{equation}
\mu_{\mathrm{ph}} = \big(\bar{\Lambda}_{\mathrm{ph}} + \beta_{\mathrm{ph},R}\tilde{\Lambda}_{q_R}\big)^{-1}\big(\bar{\Lambda}_{\mathrm{ph}} \bar{\mu}_{\mathrm{ph}} + \beta_{\mathrm{ph},R}\tilde{\Lambda}_{q_R}\Omega_{\mathrm{ph},R}\mu_R\big) \to \Omega_{\mathrm{ph},R} \mu_R,
\end{equation}
the limit holding when $\beta_{\mathrm{ph},R}\tilde{\Lambda}_{q_R} \gg \bar{\Lambda}_{\mathrm{ph}}$.
Each human observer $h_k$ subsequently equilibrates against the channel in the absorption-dominated regime $\beta_{h_k,\mathrm{ph}}\tilde{\Lambda}_{q_{\mathrm{ph}}} \gg \bar{\Lambda}_{h_k}$, giving $\mu_{h_k} \to \Omega_{h_k,\mathrm{ph}}\mu_{\mathrm{ph}}$. Composing through the channel,
\begin{equation}
\mu_{h_k} \to \Omega_{h_k,\mathrm{ph}} \Omega_{\mathrm{ph},R} \mu_R = \Omega_{h_k R} \mu_R,
\label{eq:rock_consensus}
\end{equation}
where the last equality is the transport composition rule $\Omega_{ik} = \Omega_{ij}\Omega_{jk}$, which follows from $\Omega_{ij} = e^{\phi_i}e^{-\phi_j}$ by cancellation of the middle factor $e^{-\phi_{\mathrm{ph}}}e^{\phi_{\mathrm{ph}}} = I$ (cf.\ Eq.~\eqref{eq:transport_def}). This cancellation is exact for frame-generated links, consistent with the cocycle identity of Lemma~\ref{thm:vanishing_holonomy}; loop holonomy enters only for the independent links on the interaction graph of Section~\ref{sec:discrete_regime_ii} and does not affect this equation. Two human observers $h_1, h_2$ therefore both equilibrate to the same transported rock state, each expressed in its own frame. This is the framework analog of cross-observer agreement on the rock's appearance. Classically the mechanism is precision-weighted opinion pooling~\cite{degroot1974}; its closest relative in quantum foundations is quantum Darwinism, in which objectivity arises from redundant records of a system carried by its environment, such as scattered photons~\cite{Zurek2009}.

\paragraph{Back-action in the rock limit.} The rock's update under coupling to a low-precision human contributes through the coupling term only:
\begin{equation}
\dot{\mu}_R^{(\mathrm{back})} = -\eta_q \Sigma_R \beta_{Rh} \tilde{\Lambda}_{q_h}(\mu_R - \Omega_{Rh}\mu_h) = O(\delta),
\label{eq:rock_back_mu}
\end{equation}
since $\Sigma_R = \delta I$ and the remaining factor is $O(1)$. For the covariance sector the natural gradient is $\tilde{\nabla}_{\Sigma_R}\mathcal{F} = 2\Sigma_R(\nabla_{\Sigma_R}\mathcal{F})\Sigma_R$ (Appendix~\ref{app:fisher_gaussian}). The coupling contribution
\begin{equation}
\nabla_{\Sigma_R}\mathcal{F}_{\mathrm{coupling}} = \tfrac{1}{2}\beta_{Rh}\big(\tilde{\Lambda}_{q_h} - \Lambda_R\big)
\end{equation}
produces
\begin{equation}
\begin{aligned}
\dot{\Sigma}_R^{(\mathrm{back})} &= -2\eta_q \Sigma_R \cdot \tfrac{1}{2}\beta_{Rh}\big(\tilde{\Lambda}_{q_h} - \Lambda_R\big)\cdot \Sigma_R \\
&= -\eta_q \beta_{Rh} \delta^2 \big(\tilde{\Lambda}_{q_h} - \delta^{-1}I\big) = \eta_q\beta_{Rh}\delta\big(I - \delta\tilde{\Lambda}_{q_h}\big).
\end{aligned}
\label{eq:rock_back_sigma}
\end{equation}
The two back-actions have different sizes in the rock's own geometry. The mean back-action is $O(\delta)$ in absolute terms and $O(\delta^{1/2})$ in the rock's Fisher metric, since a mean displacement $\Delta\mu$ has Fisher length $|\Delta\mu|/\delta^{1/2}$. The covariance back-action is $O(\delta)$ in absolute terms but of order one relative to $\Sigma_R$: $\dot{\Sigma}_R^{(\mathrm{back})}\Sigma_R^{-1} \to \eta_q\beta_{Rh}I$ as $\delta \to 0$. At stationarity of $\mathcal{F}_R = D_{\mathrm{KL}}(q_R\|p_R) + \beta_{Rh}D_{\mathrm{KL}}(q_R\|\Omega_{Rh}q_h)$, the covariance condition $\tfrac{1}{2}(\bar{\Lambda}_{p_R} - \Lambda_R) + \tfrac{1}{2}\beta_{Rh}(\tilde{\Lambda}_{q_h} - \Lambda_R) = 0$ gives
\begin{equation}
\Lambda_R^* = \frac{\bar{\Lambda}_{p_R} + \beta_{Rh}\tilde{\Lambda}_{q_h}}{1 + \beta_{Rh}},
\end{equation}
so with $\bar{\Lambda}_{p_R} = \delta^{-1}I$ the stationary precision is $\delta^{-1}(I + \delta\beta_{Rh}\tilde{\Lambda}_{q_h})/(1 + \beta_{Rh})$, reduced relative to the prior precision by a factor that tends to $1/(1 + \beta_{Rh})$ as $\delta \to 0$. The Fisher-Rao distance between the uncoupled and the coupled rock covariance tends to $\sqrt{K_q/2}\log(1 + \beta_{Rh})$, which does not vanish in the high-precision limit. The cause is that each relative-entropy term contains $-H[q_R]$, so coupling to any agent reweights the rock's entropy. Precision dominance therefore suppresses belief-mean back-action but not precision back-action.

\paragraph{Asymptotic reading.} The derivations above are stated in the $\delta \to 0$ rock limit. At finite $\delta$ the coupling pull on $\mu_R$ and the mean back-action are nonzero at $O(\delta)$, while the precision back-action remains of order one at every $\delta$. The claim is that this scaling sets the regime in which insensitivity to low-precision neighbors, consensus, and negligible mean back-action appear, not that these quantities vanish at finite $\delta$.

Within this interpretive extension the example supports three readings: observation is ordinary agent-agent coupling, not a special process; the observer-observed distinction is conventional, not ontological, since both are agents coupled through information exchange; and macroscopic objects are meta-agents formed from microscopic constituents through consensus. The derivation adds one framework-internal result, that precision dominance protects a high-precision agent's belief mean but not its precision.

\section{Mathematical Details}
\label{app:math}

\subsection{Variational Free Energy Preliminaries}
\label{app:variational_prelim}

Given a generative model $p(o, x)$ over observations $o$ and latent states $x$, exact inference by Bayes' rule $p(x \mid o) = p(o, x)/p(o)$ is intractable because the evidence $p(o) = \int p(o, x)\mathrm{d}x$ requires integrating over all latent states. For any approximate posterior $q(x)$ the variational free energy is an upper bound on the surprisal,
\begin{equation}
\begin{aligned}
F[q] &= \mathbb{E}_{q(x)}[\log q(x) - \log p(o, x)] \\
&= D_{\mathrm{KL}}(q(x) \| p(x \mid o)) - \log p(o) \geq -\log p(o),
\end{aligned}
\end{equation}
the inequality following from non-negativity of the relative entropy; minimizing $F$ over $q$ simultaneously tightens the evidence lower bound $\log p(o) \geq -F[q]$ and drives $D_{\mathrm{KL}}(q \| p(\cdot \mid o)) \to 0$. The measure-theoretic form of this identity, with its integrability conditions, is the classical finite evidence identity of~\cite[Ch.~5]{DennisGaugeVFE2026}. The free energy decomposes into a complexity term and an accuracy term,
\begin{equation}
F[q] = D_{\mathrm{KL}}(q(x) \| p(x)) - \mathbb{E}_{q(x)}[\log p(o \mid x)],
\end{equation}
the first penalizing departure from the prior and the second rewarding fit to observations~\cite{Friston2010, parr2022active}.

\subsection{Gaussian KL Divergence}
\label{app:gaussian_kl}

For multivariate Gaussian distributions $q = \mathcal{N}(\mu_q, \Sigma_q)$ and $p = \mathcal{N}(\mu_p, \Sigma_p)$ in $\mathbb{R}^{K_q}$, the Kullback-Leibler divergence is
\begin{equation}
D_{\mathrm{KL}}(q \| p) = \frac{1}{2}\left[\log\frac{|\Sigma_p|}{|\Sigma_q|} + \operatorname{tr}(\Sigma_p^{-1}\Sigma_q) + (\mu_p - \mu_q)^\top\Sigma_p^{-1}(\mu_p - \mu_q) - K_q\right],
\end{equation}
the Gaussian reverse relative entropy of~\cite[Ch.~G4]{DennisGaugeVFESuppG2026}. This formula is used throughout for the self and alignment terms of the consensus functional.

\subsection{Fisher Information Matrix for Gaussians}
\label{app:fisher_gaussian}

The Fisher information matrix for Gaussian distributions parametrized by $\theta = (\mu, \Sigma)$ has the block structure
\begin{equation}
F = \begin{pmatrix}
\Sigma^{-1} & 0 \\
0 & \frac{1}{2}(\Sigma^{-1} \otimes \Sigma^{-1})
\end{pmatrix},
\end{equation}
and every affine map $(\mu, \Sigma) \mapsto (Q\mu + v, Q\Sigma Q^\top)$ with $Q \in \mathrm{GL}(K_q)$ is an isometry of it~\cite[Ch.~G5, Fisher metric in the moment chart]{DennisGaugeVFESuppG2026}. The covariance block is to be read as the Fisher quadratic form restricted to symmetric perturbations, $g(d\Sigma, d\Sigma) = \tfrac{1}{2}\operatorname{tr}(\Sigma^{-1} d\Sigma \Sigma^{-1} d\Sigma)$. As an operator on symmetric matrices, $X \mapsto \tfrac{1}{2}\Sigma^{-1}X\Sigma^{-1}$ has inverse $X \mapsto 2\Sigma X \Sigma$, so the natural gradient is
\begin{align}
\tilde{\nabla}_\mu \mathcal{F} &= \Sigma (\nabla_\mu \mathcal{F}), \\
\tilde{\nabla}_\Sigma \mathcal{F} &= 2\Sigma (\nabla_\Sigma \mathcal{F}) \Sigma,
\end{align}
where $\nabla_\mu \mathcal{F}$ and $\nabla_\Sigma \mathcal{F}$ are the standard Euclidean gradients, with $\nabla_\Sigma \mathcal{F}$ the symmetric matrix gradient. For the non-redundant $\mathrm{vech}(\Sigma)$ parameterization the matrix representation of this block is $\tfrac{1}{2}\mathcal{D}^\top(\Sigma^{-1} \otimes \Sigma^{-1})\mathcal{D}$ with $\mathcal{D}$ the duplication matrix, and its matrix inverse is not simply $2(\Sigma \otimes \Sigma)$; the inverse as an operator on symmetric matrices, and hence the displayed natural gradient, is unchanged.

\subsection{SO(3) Generators and Transport Operators}

For the compact group $G = \mathrm{SO}(3)$ acting through its irreducible representation of spin $j$, the belief dimension is $K_q = 2j + 1$, and the generators $\{G_x, G_y, G_z\}$ are real skew-symmetric $K_q \times K_q$ matrices satisfying
\begin{equation}
[G_a, G_b] = \epsilon_{abc} G_c.
\end{equation}
Transport operators are computed by the matrix exponential,
\begin{equation}
\Omega_{ij} = \exp(\phi_i) \exp(-\phi_j),
\end{equation}
where $\phi_i = \phi_i^x G_x + \phi_i^y G_y + \phi_i^z G_z \in \mathfrak{so}(3)$ in the spin-$j$ representation.

The product $\exp(\phi_i)\exp(-\phi_j)$ may be replaced by the first-order Baker-Campbell-Hausdorff truncation $\exp(\phi_i - \phi_j)$, with error controlled by the leading commutator, $\Omega_{ij} = \exp\big(\phi_i - \phi_j - \tfrac{1}{2}[\phi_i, \phi_j] + \mathcal{O}(\|\phi\|^3)\big)$. The truncated transport is not of the form $U_iU_j^{-1}$: it breaks the cocycle identity, so its holonomy around a triangle of agents is of the order of $\|[\phi_i, \phi_j]\|$, and it is not invariant under the right translation $U_i \mapsto U_i g$. The vanishing field strength of Regime I (Lemma~\ref{thm:vanishing_holonomy}) holds for the exact two-exponential product only. The truncation is acceptable only when the frames are small enough that this holonomy is negligible for the purpose at hand; otherwise the exact product is used.

\subsection{Numerical Stability and Cholesky Parametrization}

Covariance matrices $\Sigma_i$ must remain positive definite throughout optimization. Two parameterizations of the covariance sector achieve this, and they differ in their dynamics. The first is the Cholesky parametrization
\begin{equation}
\Sigma_i = L_i L_i^\top,
\end{equation}
where $L_i$ is lower triangular with positive diagonal entries and gradients are taken with respect to $L_i$. Any lower-triangular $L_i$ gives a positive-semidefinite $\Sigma_i$, but $\Sigma_i$ is positive definite only when every diagonal entry of $L_i$ is nonzero, and a Euclidean gradient step on $L_i$ can drive a diagonal entry through zero. Strict positive-definiteness is therefore enforced by writing the diagonal through a diffeomorphism of the real line onto the positive half-line, such as softplus or the exponential, so that the stored degrees of freedom are unconstrained while the resulting diagonal entries cannot cross zero in exact arithmetic. The absolute value is not admissible: it vanishes at zero, is not differentiable there, and is two-to-one, so the map from stored parameters to $\mathcal{S}_{++}^{K_q}$ would not be a diffeomorphism. The second parameterization is the Fisher-Rao natural-gradient step with the affine-invariant exponential retraction of Eq.~\eqref{eq:gauge_natural_gradient} and Appendix~\ref{app:fisher_gaussian}.

With a diffeomorphic positive map on the diagonal, the map from the stored parameters to $\mathcal{S}_{++}^{K_q}$ is a diffeomorphism, so the chain-rule gradient with respect to $L_i$ vanishes exactly when $\partial\mathcal{F}/\partial\Sigma_i = 0$ and the two schemes share stationary points. They do not share trajectories, since the Euclidean step on $L_i$ induces a different metric on the covariance sector, and they do not share discrete stability. For the model free energy $\tfrac{1}{2}[\operatorname{tr}(M\Sigma) - 2\log\det\Sigma]$ with stationary point $\Sigma^* = 2M^{-1}$, the whitened linearization of the exponential-map update is multiplication by $1 - 2\eta_\Sigma$, so the update is linearly stable for $\eta_\Sigma < 1$ independently of conditioning; a Euclidean step on $\Sigma$ is stable only for $\eta_\Sigma < 2\lambda_{\min}(\Sigma^*)^2$, and a Euclidean step on $L$, with $\Sigma^*$ normalized to $\lambda_{\max}(\Sigma^*) = 1$, only for a step of order $\lambda_{\min}(\Sigma^*)$. Positive-definiteness, not step-size stability, is what the Cholesky parametrization guarantees. Under the Cholesky-Euclidean scheme the covariance relaxation rate also depends on conditioning, so a fixed ratio of learning rates does not fix the ratio of physical relaxation times. Trajectory-level natural-gradient statements therefore apply to the retraction path only; results obtained under the Cholesky-Euclidean scheme are stationarity-equivalent but follow different descent paths.

